\section{Related Work}

We organize the discussion into three parts: (i) federated learning under heterogeneous data and systems constraints, (ii) communication-efficient distributed and federated optimization, and (iii) memory-efficient training and its emerging intersection with federated learning. This structure highlights that existing methods address either communication or memory, but rarely both in a unified way under realistic federated constraints.

\subsection{Federated Learning under Heterogeneity}

Federated learning (FL) enables collaborative model training across a large population of clients without sharing raw data, by iterating local stochastic gradient steps and periodic model aggregation on a central server. The seminal FedAvg algorithm of McMahan et al.\ proposes this paradigm and demonstrates its effectiveness on mobile keyboard prediction and other on-device tasks \cite{mcmahan2017communication}, building on earlier work on communication-efficient distributed learning \cite{konecny2016federated}.

\paragraph{Data heterogeneity and client drift.}
In realistic deployments, client data are highly non-identically distributed (non-IID), which induces \emph{client drift}: local updates move in directions misaligned with the global gradient, slowing convergence and potentially harming generalization. A line of work explicitly addresses this effect. FedProx augments local objectives with a proximal term to limit local deviation from the current global model and mitigate both statistical and systems heterogeneity \cite{sahu2018fedprox}. SCAFFOLD introduces control variates maintained on both clients and server to correct the drift between local and global gradients, and proves convergence rates that match centralized SGD under strong heterogeneity \cite{karimireddy2020scaffold}. Follow-up works study the convergence of FedAvg and its variants on non-IID data, clarifying the roles of client sampling, stepsize, and local epoch length \cite{li2019convergencefedavg,yang2021linear}. Wang et al.\ further analyze objective inconsistency when different clients perform different numbers of local steps and propose adjustments to the aggregation weights to restore convergence guarantees \cite{wang2020tackling}.

\paragraph{System heterogeneity and partial participation.}
Beyond statistical heterogeneity, FL must deal with heterogeneous client compute, memory, and connectivity. The standard FedAvg protocol already assumes partial participation, where only a subset of clients join each communication round \cite{mcmahan2017communication}. Recent theory shows that linear speedup in the number of clients is still achievable under non-IID data and partial participation if the sampling and stepsizes are chosen appropriately \cite{yang2021linear}. System-level works design algorithms that are robust to stragglers and varying computation budgets, for example by adjusting local epoch numbers, using proximal or dynamic regularization \cite{sahu2018fedprox,acar2021feddyn}, or adopting primal--dual formulations such as FedADMM that explicitly incorporate partial participation and communication constraints in the optimization problem \cite{wang2022fedadmm}. Zheng et al.\ analyze consensus-based FL over heterogeneous data, providing a complementary convergence perspective in terms of consensus error across clients \cite{zheng2023consensusFL}.

\paragraph{Personalized and hierarchical federated learning.}
When client data distributions differ significantly, a single global model can underperform for many users. Personalized FL aims to balance shared global structure with client-specific adaptation. Per-FedAvg and related meta-learning approaches optimize a global initialization that is quickly adaptable to each client via a few gradient steps \cite{zhang2020perfedavg}. Dinh et al.\ formulate personalization via Moreau envelopes, leading to an algorithm (pFedMe) that optimizes client-specific models regularized toward a shared global model and enjoys convergence guarantees \cite{dinh2020pfedme}. Layer-wise and representation-level aggregation techniques further disentangle global and local components \cite{ma2022layerwise}, while recent works exploit distillation-based personalization and clustering to handle system and model heterogeneity \cite{chen2024spectralcodistill,xie2024MHpFLID}. These personalization mechanisms, however, typically introduce additional model copies, auxiliary heads, or teacher--student structures, increasing the memory footprint on each client.

\paragraph{Limitations with respect to memory.}
Overall, the FL literature has focused primarily on handling data and systems heterogeneity while preserving convergence rates and sometimes improving communication efficiency \cite{mcmahan2017communication,li2019convergencefedavg,yang2021linear,karimireddy2020scaffold,acar2021feddyn}. Much less attention has been paid to the \emph{memory} footprint of local training---including activations, gradients, and optimizer states---which is a critical bottleneck for low-end devices and for training increasingly large models (e.g., foundation models) in federated settings. Existing FL algorithms typically assume that clients can store full-precision model parameters, gradients, and optimizer states, and they do not provide a principled framework for trading off memory against communication and convergence.

\subsection{Communication-Efficient Distributed and Federated Optimization}

Communication is a primary bottleneck in both data-center and cross-device distributed training. Two general strategies have emerged: reducing the \emph{frequency} of communication via local updates, and reducing the \emph{per-round} communication cost via compression or sparsification.

\paragraph{Local SGD and communication frequency.}
Local SGD performs multiple local updates between synchronization steps, thereby amortizing communication over several gradient computations \cite{stich2019localSGD}. Stich shows that, under appropriate choices of local steps and stepsizes, local SGD can achieve nearly the same convergence rate as synchronous mini-batch SGD while using significantly fewer communication rounds \cite{stich2019localSGD}. Gao et al.\ extend the analysis of local SGD to the FL setting with non-IID data, characterizing how communication-efficient local updates interact with client heterogeneity \cite{gao2021localSGDfed}. SlowMo combines local SGD with a slow momentum update on the server side, improving robustness to stale or noisy local updates \cite{wang2019slowmo}. Qsparse-Local-SGD further unifies local computation with gradient sparsification and quantization, showcasing the benefits of combining multiple communication-saving techniques in a single framework \cite{basu2019qsparselocal}.

\paragraph{Gradient compression and error feedback.}
A complementary line of work compresses each communicated update. QSGD proposes randomized quantization schemes that provide unbiased gradient estimates with controllable variance, enabling theoretical convergence guarantees for quantized SGD \cite{alistarh2017qsgd}. Sketching-based methods further reduce communication by sending compact sketches of large gradients, with recovery guarantees for the largest coordinates \cite{li2019sketchedsgd}. However, many practical compressors---including top-$k$ sparsification and sign-based quantization---are biased or contractive. Error-feedback (EF) schemes correct this bias by accumulating the compression error and feeding it back into subsequent updates; EF has been shown to restore convergence guarantees even for heavily compressed gradients \cite{karimireddy2019errorfeedback}. Recent refinements such as EF21 simplify the mechanism and improve both theory and practice \cite{richtarik2021ef21}. Momentum-based compressed methods, such as blockwise momentum SGD with error-feedback \cite{zheng2019blockmomentumEF}, further enhance convergence in heterogeneous settings.

\paragraph{Decentralized and topology-aware training.}
Decentralized optimization eliminates the central parameter server and replaces it with peer-to-peer communication over a network graph. D-PSGD and its asynchronous variants show that decentralized SGD can outperform centralized schemes in heterogeneous environments by avoiding server bottlenecks and distributing communication more evenly \cite{lian2017dpsgd}. Nevertheless, early analyses revealed that data heterogeneity could severely degrade convergence in decentralized training. D$^2$ (Decentralized Training over Decentralized Data) addresses this by carefully designing gradient tracking and update rules that remove the influence of data heterogeneity from the error bounds \cite{tang2018d2}. A rich literature on compressed and variance-reduced decentralized methods further optimizes communication by combining gradient tracking with contractive compressors \cite{li2019communicationGradientTracking,liao2022compressedDecentralized}.

More recently, dynamic-topology approaches aim to scale decentralized learning to very large networks. Teleportation periodically activates only a small subset of edges while preserving fast convergence through random long-range connections, achieving scalable decentralized learning with low communication \cite{takezawa2025teleportation}. Other works jointly optimize collaboration patterns and communication topology in decentralized FL, dynamically selecting neighbors or clusters based on utility criteria \cite{jiang2025dfedmq}.

\paragraph{Lower bounds and near-optimal algorithms.}
Theoretical lower bounds for distributed optimization with communication compression clarify fundamental trade-offs between compression levels, network size, and convergence rates. Huang et al.\ derive lower bounds on the communication complexity of distributed learning under compressed communication and propose nearly optimal algorithms that match these bounds up to logarithmic factors \cite{he2022lowerbounds}. Follow-up work extends these results to stochastic and bilevel optimization, further refining our understanding of when and how unbiased or contractive compression can save communication without sacrificing optimal rates \cite{he2023lowerboundsStoch,he2024distributedBilevel}. These results motivate algorithm designs such as NEOLITHIC-style methods that approach the theoretical limits by carefully integrating compression, variance reduction, and gradient tracking.

\paragraph{Limitations with respect to memory.}
Despite substantial progress, most communication-efficient methods assume that local gradient computation and optimizer state storage are unconstrained. Clients must still materialize full gradients and maintain full optimizer states (e.g., momentum or adaptive moments) in memory. Even when low-rank or sketch-based compression is used for communication \cite{alistarh2017qsgd,li2019sketchedsgd,basu2019qsparselocal}, the compression is typically applied \emph{after} full gradients are computed and stored. Thus, these methods do not directly address the memory overhead of training large models on resource-limited devices, nor do they provide convergence guarantees under joint memory and communication constraints.

\subsection{Memory-Efficient Training and Federated Learning}

Orthogonal to communication efficiency, a separate line of work aims to reduce the memory footprint of training deep networks. These methods target activations, model parameters, and optimizer states, and only recently have begun to intersect with federated learning.

\paragraph{Sublinear-memory and model-compression techniques.}
Classical techniques such as pruning, quantization, and parameter sharing compress model parameters for deployment, but typically assume unconstrained memory during training. Chen et al.\ show that training deep networks can be implemented with $\mathcal{O}(\sqrt{n})$ memory in the number of layers $n$ via gradient checkpointing and recomputation, at the cost of extra forward passes per mini-batch \cite{chen2016sublinear}. Hash-based parameter sharing further reduces model memory while preserving accuracy by hashing multiple weights into shared parameters \cite{chen2015hashing}. These algorithmic ideas are widely used in practical systems, but they do not directly account for federated constraints such as partial participation and heterogeneous devices.

\paragraph{Low-rank adaptation and optimizer-state reduction.}
Parameter-efficient fine-tuning (PEFT) approaches such as LoRA freeze the pretrained backbone and insert low-rank adapters, dramatically reducing the number of trainable parameters and the optimizer state memory required for adaptation \cite{hu2021lora}. Building on this idea at the optimization level, GaLore proposes to project gradients onto a low-rank subspace before applying updates, so that only low-rank optimizer states need to be stored \cite{zhao2024galore}. This leads to substantial memory savings in the training of large language models (LLMs), with minimal impact on convergence in practice. However, subsequent work reveals that naive subspace optimization can fail to converge under stochastic noise and proposes variants with provable convergence guarantees by carefully controlling the subspace dynamics \cite{he2024subspaceLLM}. These developments connect subspace methods to classical geometric convergence theory for preconditioned steepest descent \cite{neymeyr2011psd}, highlighting that low-rank constraints on the optimization trajectory must be designed with care.

\paragraph{Memory-efficient federated learning with low-rank and split training.}
Combining low-rank adaptation with FL is a natural way to reduce both communication and memory costs on clients. Recent works explore training only low-rank adapters on each client while freezing the backbone, thereby shrinking the transmitted updates and the local optimizer states. FLoCoRA studies the application of LoRA-style adapters in FL from scratch, analyzing how low-rank parameterization interacts with client heterogeneity and communication constraints \cite{ribeiro2024flocora}. Trautmann et al.\ investigate how to aggregate low-rank adapters across clients for federated fine-tuning of language models, proposing aggregation strategies that respect the low-rank structure \cite{trautmann2024aggregateLora}. In parallel, Chen et al.\ propose a memory-efficient split FL framework for LLM fine-tuning on heterogeneous mobile devices, where low-rank adaptation is combined with model splitting between devices and edge servers to alleviate memory pressure on both sides \cite{chen2025memsplitfl}.

While these methods demonstrate the practicality of low-rank or split training in federated settings, their theoretical understanding remains limited. In particular, when local updates are restricted to low-dimensional subspaces or involve non-linear transformations (e.g., orthogonalization or reparameterization), the standard drift-correction analyses used in FL (as in SCAFFOLD \cite{karimireddy2020scaffold}) no longer apply directly. Existing works typically either rely on empirical evaluation or perform analysis in simplified centralized settings \cite{zhao2024galore,he2024subspaceLLM}, leaving open the question of how to design drift-correction mechanisms tailored to memory-efficient local updates in FL.

\paragraph{Personalization, activation memory, and dynamic participation.}
Memory constraints also interact with personalization and activation storage. Distillation-based personalized FL often requires maintaining teacher and student models, or auxiliary classifiers, which can significantly increase per-client memory usage \cite{chen2024spectralcodistill,xie2024MHpFLID}. Activation checkpointing \cite{chen2016sublinear} can be used on clients to reduce activation memory, but the interplay between checkpointing schedules, partial participation, and non-IID data remains largely unexplored. Dynamic-topology methods such as Teleportation can reduce the number of active communication links per round \cite{takezawa2025teleportation}, and topology-optimization approaches in decentralized FL can limit the number of neighbors each client interacts with \cite{jiang2025dfedmq}, implicitly reducing the memory needed for storing neighbor states. However, these system-level techniques do not directly reduce the core memory footprint of model parameters, gradients, and optimizer states.

\paragraph{Open gaps and motivation.}
In summary, the literature provides sophisticated tools for handling heterogeneity and reducing communication in FL \cite{mcmahan2017communication,karimireddy2020scaffold,stich2019localSGD,alistarh2017qsgd,he2022lowerbounds}, and a rapidly growing body of work on memory-efficient training via low-rank subspace methods and checkpointing in centralized settings \cite{chen2016sublinear,hu2021lora,zhao2024galore,he2024subspaceLLM}. However, there is still no unified framework that:
(i) jointly optimizes memory and communication in FL;
(ii) provides principled drift-correction mechanisms for memory-efficient local updates under non-IID data and partial participation; and
(iii) achieves provable convergence under realistic nonconvex objectives and heterogeneous systems constraints.

Addressing these gaps requires new algorithmic designs that couple subspace or low-rank updates with compression and local training, alongside new analyses that extend control-variate and gradient-tracking ideas to non-Euclidean, memory-constrained local optimization steps.
